% GraphHunter — Libro «Algoritmos»
% Textbook pedagógico estilo CLRS: del concepto básico al algoritmo
% avanzado, con teoremas, lemas, ejemplos trabajados, ejercicios y
% soluciones selectas en apéndice.
%
% Color de acento: verde bosque #1F5D3A (distinto a los otros 5 libros).

\usepackage{fontspec}
\setmainfont{Latin Modern Roman}
\setsansfont{Latin Modern Sans}
\setmonofont{Latin Modern Mono}
\usepackage[spanish,es-tabla,es-lcroman,es-noshorthands]{babel}
\usepackage{csquotes}

\usepackage[a4paper,margin=2.4cm]{geometry}

\usepackage{amsmath, amssymb, amsthm}
\usepackage{mathtools}

\usepackage[ruled,linesnumbered,algochapter]{algorithm2e}
\SetKwComment{Comment}{$\triangleright$ }{}
\SetKwFor{For}{for}{do}{end}
\SetKwFor{While}{while}{do}{end}
\SetKwIF{If}{ElseIf}{Else}{if}{then}{else if}{else}{end}
\SetKwRepeat{Repeat}{repeat}{until}

\usepackage{listings}
\usepackage{xcolor}

% Acento «verde bosque» — distinto a:
%   diseño    (#4B2E83 morado),
%   paper / arquitectura / spec / manual (azules y rojos varios)
\definecolor{accent}{HTML}{1F5D3A}
\definecolor{accentlight}{HTML}{E0EFE5}
\definecolor{ejemplostroke}{HTML}{8AB89A}
\definecolor{exercisestroke}{HTML}{1F5D3A}
\definecolor{codecomment}{rgb}{0.35,0.45,0.45}
\definecolor{codekw}{rgb}{0.10,0.40,0.20}
\definecolor{codestr}{rgb}{0.55,0.10,0.10}

\lstdefinelanguage{Rust}{
    morekeywords={fn,let,mut,pub,struct,enum,impl,trait,where,if,else,for,while,
                  loop,match,return,use,mod,const,static,self,Self,as,in,ref,move,
                  async,await,dyn,unsafe,crate,super,extern,type,break,continue},
    morecomment=[l]{//},
    morecomment=[s]{/*}{*/},
    morestring=[b]",
    morestring=[b]',
}
\lstset{
    basicstyle=\ttfamily\small,
    keywordstyle=\color{codekw}\bfseries,
    commentstyle=\color{codecomment}\itshape,
    stringstyle=\color{codestr},
    showstringspaces=false,
    frame=single,
    numbers=none,
    captionpos=b,
    breaklines=true,
    tabsize=4,
    extendedchars=true,
    literate=%
        {α}{{$\alpha$}}1
        {β}{{$\beta$}}1
        {ρ}{{$\rho$}}1
        {Σ}{{$\Sigma$}}1
        {≤}{{$\leq$}}1
        {≥}{{$\geq$}}1
        {→}{{$\rightarrow$}}1
        {↔}{{$\leftrightarrow$}}1,
}

\usepackage{hyperref}
\hypersetup{
    colorlinks=true,
    linkcolor=accent,
    citecolor=black,
    urlcolor=accent,
    pdftitle={GraphHunter — Algoritmos},
    pdfauthor={Lucas Sotomayor / Base4Security},
    bookmarks=true,
    bookmarksopen=true,
    bookmarksnumbered=true,
}
\usepackage{bookmark}

\usepackage{booktabs}
\usepackage{longtable}
\usepackage{float}
\usepackage{tikz}
\usetikzlibrary{arrows.meta, positioning, shapes.geometric, fit, backgrounds, calc, shapes.misc, decorations.pathreplacing}

% --- Cajas de callout -------------------------------------------------
\usepackage[most]{tcolorbox}

% «Intuición»: resumen plain-language antes de la formalización.
\newtcolorbox{intuicion}{
    colback=accentlight,
    colframe=accent,
    arc=2mm,
    boxrule=0.6pt,
    left=4pt, right=4pt, top=4pt, bottom=4pt,
    fonttitle=\bfseries,
    title=Intuición
}

% «Ejemplo»: ejemplo trabajado extenso. Estilo distinto a Intuición:
% fondo blanco, borde sólido, para diferenciar visualmente. El
% argumento opcional pone un sufijo al título: \begin{ejemplo}[K_4
% como join] genera "Ejemplo: K_4 como join".
\NewTColorBox{ejemplo}{ !O{} }{
    colback=white,
    colframe=ejemplostroke,
    arc=1mm,
    boxrule=0.6pt,
    left=4pt, right=4pt, top=4pt, bottom=4pt,
    fonttitle=\bfseries,
    title={Ejemplo\IfValueTF{#1}{\ifx\relax#1\relax\else: #1\fi}{}}
}

% «Ejercicios»: bloque al final de cada capítulo. Llevará una lista
% \begin{enumerate} adentro con \item enunciados. Las soluciones a los
% ejercicios marcados con \starred{} viven en el Apéndice C.
\newtcolorbox{ejercicios}{
    enhanced,
    colback=white,
    colframe=exercisestroke,
    arc=1mm,
    boxrule=0.8pt,
    left=4pt, right=4pt, top=4pt, bottom=4pt,
    fonttitle=\bfseries,
    title=Ejercicios,
    breakable
}

% «Solución»: usado en el Apéndice C. Cada solución va dentro.
\newenvironment{solucion}[1]{%
    \par\vspace{0.5em}\noindent
    \textbf{Solución \thechapter.#1.}\enspace\ignorespaces
}{\par\vspace{0.5em}}

% --- Teoremas, lemas, definiciones ------------------------------------
\theoremstyle{plain}
\newtheorem{theorem}{Teorema}[chapter]
\newtheorem{lemma}[theorem]{Lema}
\newtheorem{corollary}[theorem]{Corolario}
\newtheorem{proposition}[theorem]{Proposición}

\theoremstyle{definition}
\newtheorem{definition}[theorem]{Definición}
\newtheorem{example}[theorem]{Ejemplo}

\theoremstyle{remark}
\newtheorem*{invariant}{Invariante}
\newtheorem*{nota}{Nota}
\newtheorem*{advertencia}{Advertencia}

% --- Macros varias ----------------------------------------------------
\newcommand{\filepath}[1]{\texttt{\small\detokenize{#1}}}
\newcommand{\code}[1]{\texttt{\detokenize{#1}}}

% Marca un ejercicio cuya solución está en el Apéndice C.
\newcommand{\starred}{$\bigstar$}

% Bloque «En el código»: parrafito que apunta del concepto al fichero
% real del repo donde vive su implementación.
\newenvironment{enelcodigo}{%
    \par\vspace{0.4em}%
    \noindent\textbf{En el código.}\nopagebreak\par\nopagebreak\vspace{0.2em}%
}{\par\vspace{0.4em}}

% Estilos de capítulo y parts.
\usepackage{titlesec}
\titleformat{\chapter}[display]
  {\normalfont\huge\bfseries\color{accent}}
  {\chaptertitlename~\thechapter}
  {20pt}
  {\Huge}
\titleformat{\part}[display]
  {\normalfont\Huge\bfseries\color{accent}\filcenter}
  {\partname~\thepart}
  {20pt}
  {\Huge}

% TODO marker (rojo, vistoso, no hay que hacer build especial).
\newcommand{\TODO}[1]{\textbf{\textcolor{red}{[TODO: #1]}}}
\newcommand{\passthrough}[1]{#1}

% Símbolos matemáticos que el libro usa con frecuencia.
\newcommand{\E}{\mathcal{E}}        % conjunto de aristas
\newcommand{\V}{\mathcal{V}}        % conjunto de vértices
\newcommand{\T}{\mathcal{T}}        % dominio temporal
\newcommand{\Lc}{\mathcal{L}}       % conjunto de etiquetas
\newcommand{\NLF}{\mathrm{NLF}}
\newcommand{\bits}{\mathrm{bits}}
\newcommand{\AGM}{\mathrm{AGM}}
\newcommand{\rhostar}{\rho^{*}}
\newcommand{\indeg}{\mathrm{in\text{-}deg}}
\newcommand{\outdeg}{\mathrm{out\text{-}deg}}
